\chapter{绪论}

\section{研究背景}
医学图像检测、分割与量化分析是医学人工智能研究中的核心内容。对于病理图像而言，细胞核检测与实例分割直接关系到组织结构理解、细胞计数、肿瘤微环境分析和后续临床量化建模；对于放射与病理融合分析而言，可靠的区域定位和结构提取又是疗效评估、风险分层和预后预测的重要基础。可以说，前端识别能力的稳定性在很大程度上决定了后续医学分析任务的上限。

然而，这一领域长期存在两个互相交织的难题。其一是监督信号昂贵且稀缺。医学图像尤其是病理图像中的实例级标注，需要专业病理医师或经过长期训练的研究人员完成，时间成本与经济成本都很高。其二是目标域迁移困难。同一疾病在不同中心、不同切片制备流程、不同染色和不同扫描仪条件下，会表现出明显不同的数据分布；即使模型在单一公开数据集上表现良好，一旦进入真实场景，也常常面临性能下降的问题。

近年来，大规模视觉语言预训练模型的兴起为解决上述问题提供了新的可能。一方面，这类模型具备一定的开放词汇识别能力，使得“未见类别识别”和“少样本适配”不再完全依赖人工标注；另一方面，提示学习机制使模型有机会在结构保持相对稳定的前提下，通过较少附加参数实现目标域迁移。这些新进展为医学图像分析带来了新的方法论转向，即从传统强监督闭集学习逐步走向弱监督、无监督和开放场景建模。

\section{研究动机}
本文的研究动机来源于一个持续演化的认识：虽然医学图像强监督标注昂贵，但目标结构、跨模态知识以及少量示例本身蕴含着可被模型利用的先验。换言之，昂贵人工标注并不是唯一的监督来源，合理设计的伪标签、自训练、文本提示、图像提示和结构先验都可能成为替代性学习信号。

沿着这一思路，本文将研究问题组织为一条从“识别入口构建”走向“临床量化解释”的主线：首先，在缺乏目标域检测标注时，如何借助视觉语言大模型的开放词汇能力与自动提示构造建立零样本细胞核检测入口？其次，当零样本病理检测对手工提示工程高度敏感、且密集细胞核易出现漏检与重叠时，如何通过自动提示机制生成更稳定的语义入口，并结合自训练迭代实现目标域适配？再次，当纯文本提示难以稳定表达细粒度医学视觉差异时，如何进一步让图像提示以兼容形式进入视觉语言模型，以增强跨域感知与推理能力？最后，当前端识别能力逐步建立后，如何将其组织为可解释、可统计的数字病理指标，用于疗效评估与预后分析？

\section{研究内容}
围绕上述主线，本文主要开展以下四方面研究。

第一，研究视觉语言大模型在细胞核检测中的迁移问题，提出基于自动属性提示与自训练的零样本细胞核检测方法。该工作解决的是“缺乏目标域检测标注时如何构造可辨识的医学提示，并将开放词汇能力转化为目标域可用的检测性能”。

第二，针对零样本检测对提示工程的敏感性以及密集细胞核场景下的漏检、重叠问题，提出自动提示框架 AttriPrompter：从无标注病理图像中自动生成形状/颜色属性，并通过同义与“程度”增强扩展描述粒度，进一步按图文相关性对提示序列排序以获得更可靠的初始定位；在此基础上，引入自训练知识蒸馏迭代优化，以伪标签驱动目标域适配并提升无监督检测性能。该工作解决的是“如何在无需人工提示模板与检测标注的条件下，自动构造高质量提示并稳定提升零样本检测效果”。

第三，研究图像提示与文本提示的统一建模与注入机制，提出视觉文本化方法，通过将少量支持图像映射到语言空间兼容的提示表示，增强模型对细粒度差异与跨域风格变化的感知能力。该工作解决的是“当文本描述不足时，如何让图像示例以提示的形式直接参与推理”。

第四，将上述方法体系进一步应用于肺鳞癌新辅助免疫化疗后的数字病理分析，构建面向残余肿瘤负荷与肿瘤周围免疫细胞密度的量化框架，并通过统计分析探索其与疗效及预后之间的关系。该工作解决的是“如何将检测、分割结果转化为临床可解释指标并形成分析闭环”。

\section{技术路线}
本文的技术路线可以概括为“从提示驱动的识别能力构建到临床导向的可解释量化”的逐步推进：前期工作面向无检测标注场景，重点解决自动提示构造与目标域自适配问题；中期工作进一步面向未见类别与跨域分布偏移，重点解决结构先验注入与提示融合问题；后期工作将上述识别能力整合到临床导向的数字病理链条中，强调可解释指标构建与统计验证。

从建模角度看，本文的方法体系有三个共同特点。第一，强调尽量减少对昂贵强监督标注的依赖。第二，强调最大化利用视觉语言预训练带来的跨模态先验，并通过提示机制进行轻量适配。第三，强调输出结果的可解释性，使方法不仅服务于 benchmark 指标，也服务于真实数字病理分析需求。

\section{论文结构安排}
全文共分七章。第一章为绪论，介绍研究背景、动机、内容和结构安排。第二章总结相关研究基础，包括视觉语言预训练、提示学习、医学检测与分割以及数字病理分析。第三章介绍自动属性提示与自训练联合的零样本细胞核检测方法。第四章介绍自动提示框架 AttriPrompter 及其自训练知识蒸馏的迭代优化方法。第五章介绍视觉文本化框架及其图像提示驱动的检测与分割方法。第六章介绍面向肺鳞癌疗效评估的数字病理量化分析框架。第七章对全文进行总结并讨论后续研究方向。

\section{主要贡献}
本文的主要贡献体现在以下四个方面。其一，提出了基于自动属性提示与自训练的零样本细胞核检测方法，为医学无标注检测提供了可扩展的实现路径。其二，提出了自动提示框架 AttriPrompter，通过“属性生成—增强—排序”自动构造高质量提示序列，并结合自训练知识蒸馏提升密集细胞核场景下的零样本检测稳定性与目标域适配能力。其三，提出了视觉文本化方法，将图像提示映射到语言空间，从而增强模型在跨域与细粒度目标场景中的感知能力。其四，构建了面向肺鳞癌疗效评估的数字病理分析框架，将前端识别结果进一步组织为具有临床解释意义的量化指标并开展统计分析。
